\clearpage\thispagestyle{empty}
\pdfbookmark[0]{Quellen zur Besprechung von Zachs Tafeln}{zach1808tables-sources}
\begin{center}\textbf{Quellen zur Besprechung von Zachs Tafeln}\end{center}
\noindent Diese KI-Übertragung ins heutige Deutsch folgt der Besprechung von Gauß auf S.~489--492 der \href{https://gdz.sub.uni-goettingen.de/id/PPN236007467}{\textit{Werke}, Bd.~VI} (Göttingen, 1874). Sie erschien ohne Namensnennung unter der Überschrift \emph{Gotha} in den \textit{Göttingischen gelehrten Anzeigen}, Stück~2, 2. Januar 1808, S.~\href{https://api.digitale-sammlungen.de/iiif/image/v2/bsb10538673_00024/full/1800,/0/default.jpg}{18}--\href{https://api.digitale-sammlungen.de/iiif/image/v2/bsb10538673_00030/full/1800,/0/default.jpg}{24}. Das verlinkte Exemplar stammt aus der Bayerischen Staatsbibliothek.
\par\smallskip\noindent Der lateinische Titel lautet auf Deutsch \emph{Spezielle Tafeln der Aberration und Nutation in Rektaszension und Deklination zur Berechnung der scheinbaren oder wahren Positionen der Fixsterne, zusammen mit einem neuen Verzeichnis von 494 der bedeutenderen Tierkreissterne, erstellt an der Ernestinischen Sternwarte für den Beginn des Jahres 1800, und weiteren einschlägigen Tafeln}, von Franz Xaver Freiherr von Zach. Die Besprechung beschreibt die Tafeln und Verzeichnisse, ohne sie abzudrucken. Im ersten Absatz steht die lateinische Wendung \emph{monumentum aere perennius}. KI-Übertragung und KI-Satz.
\clearpage
